\section{Aula 7}

Serei um pouco mais breve nessas notas, pois essa aula
foi mais expositiva. Vimos percolação em redes hierárquicas.
Começamos com um grafo $G=G_1$ conexo com dois vértices especiais
$a$ e $b$. Definimos um grafo $G_i$ recursivamente
(considerando uma orientação de $G$  e suas arestas)
trocando todas as arestas ${u,v}$ de $G_{i-1}$ por uma cópia de $G$
onde (podendo mudar com orientação) $u$ se torna $a$ e $v$ se
torna $b$.

O nosso exemplo foi o diamante, Figura 1 de \url{https://arxiv.org/pdf/2606.11503}.

O nosso evento de interesse agora é conectividade entre os vértices
$a_i,b_i$ quando
percolamos arestas independentemente com probabilidade $p$. Mais precisamente,
queremos entender bem
\[
	f_i(p) = \Pr_p^i(a_i \leftrightarrow b_i).
\]
É fácil ver que uma dinâmica é estabelecida nessas probabilidades. Se pensarmos
que $G_{i}$ é uma cópia de $G$ em que trocamos cada aresta por $G_{i-1}$,
sacamos imediatamente que $f_i(p) = f_1(f_{i-1}(p))$. Para facilitar
notação, chamaremos $f_1$ somente de $f$. Isso significa que estamos
interessados
precisamente nas iteradas de $f$ e queremos entender esse sistema dinâmico.

O Augusto, o Gugu e seus coautores mostraram que se o grafo satisfaz
duas condições básicas:
\begin{itemize}
	\item  $d(a,b) \geq 2$.
	\item O corte mínimo entre $a$ e $b$ usa mais de uma aresta.
\end{itemize}
então existem exatamente três pontos fixos $0, 1$ e $p_c$. Onde $0$ e
$1$ são atratores
e $p_c^*$ é hiperbólico. É fácil ver pela condição $1$ que se
estivermos suficientemente
próximos de $0$, somos atraídos para $0$ em velocidade duplo
exponencial. Mais precisamente,
perto de $0$, $f(p) \leq Cp^2$, donde segue que se $p < 1/C$, então
$f^{k}(p) \leq \exp(-\alpha2^k)$. Isso pode parecer rápido demais, mas o Augusto
fez o comentário que eu considero mais importante dessa aula. Nesse
modelo, o número de vizinhos
de $a_k$ é da ordem de $2^k$. De fato, é como se estivessemos olhando
para um quadrado em $\Z^2$
e estamos interssados na probabilidade de cruzar esse quadrado, onde
$a_k$ e $b_k$ representam
lados opostos. Vendo dessa forma, devemos realmente pensar que $2^k$
faz o papel de $n$
em percolação normal e estamos obtendo aqui um decaimento exponencial.

Esse sketch já serve para qualquer $p < p_c$, pois a função é
crescente e admite somente
$0$ como ponto crítico. Então para todo ponto, em tempo finito
conseguimos esse decaimento.
Exatamente a mesma cois a pode ser observada para $p > p_c$ e
considerando a função $1 - (1-p)^2$.

Depois vimos estratégias usadas para calcular mais precisamente os
expoentes críticos, nada foi muito
formal mas é bem bonitinho. Eles foram bem espertinhos em como fazer
isso. Escolhendo uma aresta $e$
aleatória em $G_k$, ficamos interessados na probabilidade de $E$
conectar com $a_k$ e $b_k$.
Definindo $X_k^a$ sendo o número de arestas que só conectam com $a$,
$X_k^b$ as que só conectam com $b$
e $X_k^{a,b}$ as que conectam com ambos. Conseguimos escrever um
sistema dinâmico tridimensional
com $(\Ex X_k^a, \Ex X_k^b, \Ex X_k^{a,b})$, estudando propriedades
da matriz de transição, dá
para checar que ela é uma soma de subestocásticas e que sua norma
cresce exponencialmente com
uma função de $p$, $||T_{k}|| = (Cp)^k$ também.

Essa aula foi divertida, vai que um dia veremos uma sobre redes de
hipergrafos \Smiley.
